\clearpage
\VerdierPhysicalPageAnchor{0020}
茎は捩れのないアーベル群である（\emph{平坦}分解）。次の事実は容易に
確かめられる：

\VerdierSemanticAnchor{intro-1-flat-common-refinement}
\noindent\emph{a)} 二つの平坦分解
$L_i^{\bullet}\xrightarrow{\varepsilon_i}G$、$i=1,2$ が与えられると、
第三の平坦分解 $L_3^{\bullet}\xrightarrow{\varepsilon_3}G$ と、
ホモトピーを除いて可換な分解の図式が存在する：

\VerdierSemanticAnchor{intro-1-flat-common-refinement-diagram}
\[
\begin{tikzcd}[column sep=large,row sep=large]
  & L_1^{\bullet} \arrow[dr,"\varepsilon_1"] & \\
L_3^{\bullet} \arrow[ur,"\rho_1"] \arrow[rr,"\varepsilon_3"]
  \arrow[dr,"\rho_2"'] && G \\
  & L_2^{\bullet} \arrow[ur,"\varepsilon_2"'] &
\end{tikzcd}
\qquad\text{.}
\]

\VerdierSemanticAnchor{intro-1-flat-homotopy-equalizer}
\noindent\emph{b)} 平坦分解の二つの射

\VerdierSemanticAnchor{intro-1-flat-parallel-morphisms-diagram}
\[
\begin{tikzcd}[column sep=large,row sep=large]
L_1^{\bullet}
  \arrow[d,bend right=12,"\rho"']
  \arrow[d,bend left=12,"\rho'"]
  \arrow[dr,"\varepsilon_1"] & \\
L_2^{\bullet} \arrow[r,"\varepsilon_2"'] & G
\end{tikzcd}
\qquad\text{,}
\]
が与えられると、第三の平坦分解

\VerdierSemanticAnchor{intro-1-flat-third-resolution}
\[
L_3^{\bullet}\xrightarrow{\varepsilon_3}G
\qquad\text{,}
\]
と、分解の射

\VerdierSemanticAnchor{intro-1-flat-sigma-diagram}
\[
\begin{tikzcd}[column sep=large,row sep=large]
L_3^{\bullet} \arrow[d,"\sigma"'] \arrow[dr,"\varepsilon_3"] & \\
L_1^{\bullet} \arrow[r,"\varepsilon_1"'] & G
\end{tikzcd}
\qquad\text{,}
\]
が存在し、複体の射 $\rho\sigma$ と $\rho'\sigma$ はホモトピックになる。

\VerdierSemanticAnchor{intro-1-flat-morphism}
\noindent\emph{c)} 任意の平坦分解の射

\VerdierSemanticAnchor{intro-1-flat-morphism-diagram}
\[
\begin{tikzcd}[column sep=large,row sep=large]
L_1^{\bullet} \arrow[d,"\rho"'] \arrow[dr,"\varepsilon_1"] & \\
L_2^{\bullet} \arrow[r,"\varepsilon_2"'] & G
\end{tikzcd}
\qquad\text{,}
\]
